%% Question bank chapter file - generated by bank_tools.py build (do not edit while the bank is still being built:
%% it is regenerated from the per-test data in Question Bank\data). Every question is one bankq block:
%%   \begin{bankq}[id=..., source={...}, paper=.., qno=.., type=..., subtopic={...}, answer={...}, answer_src=key|solved,
%%                 working={...}, difficulty=1-10, diff_parts={K S I C F}, diff_note={...},
%%                 group=<passage id or empty>, also_in={duplicate copies in other tests}]
%% difficulty = min(10, K+S+I+C+F-1): K concepts 1-3, S steps 1-4, I insight 0-2, C computation 0-2, F format 0-1
%%   ...question LaTeX, options, figures (\bfig / \bfigi -> img/)...
%%   \end{bankq}
%% Passages shared by several questions are bankpassage blocks; their questions carry group=<passage id>.
\documentclass[10pt]{article}
\usepackage[a4paper,margin=1.8cm]{geometry}
\usepackage[T1]{fontenc}
\usepackage{amsmath,amssymb,mathtools}
\usepackage{graphicx}
\usepackage[export]{adjustbox}
\usepackage{xcolor,needspace,keyval}
\usepackage{enumitem}
\usepackage{booktabs,array,tabularx}
\usepackage[hidelinks]{hyperref}
\setlength{\parindent}{0pt}
\setlength{\parskip}{3pt}
\makeatletter
\define@key{bq}{id}{\def\bq@id{#1}}\define@key{bq}{source}{\def\bq@src{#1}}\define@key{bq}{paper}{\def\bq@paper{#1}}
\define@key{bq}{qno}{\def\bq@qno{#1}}\define@key{bq}{type}{\def\bq@type{#1}}\define@key{bq}{subtopic}{\def\bq@sub{#1}}
\define@key{bq}{answer}{\def\bq@ans{#1}}\define@key{bq}{answer_src}{\def\bq@asrc{#1}}\define@key{bq}{working}{\def\bq@work{#1}}
\define@key{bq}{group}{\def\bq@group{#1}}\define@key{bq}{also_in}{\def\bq@also{#1}}
\define@key{bq}{difficulty}{\def\bq@diff{#1}}\define@key{bq}{diff_parts}{\def\bq@dparts{#1}}\define@key{bq}{diff_note}{\def\bq@dnote{#1}}
\newenvironment{bankq}[1][]{\def\bq@work{}\def\bq@also{}\def\bq@group{}\def\bq@diff{?}\def\bq@dparts{}\def\bq@dnote{}\setkeys{bq}{#1}%
  \par\needspace{4\baselineskip}\medskip\noindent{\small\textbf{\bq@id}\hfill\color{gray}[\bq@src, Paper \bq@paper, Q\bq@qno]
  \bq@type\ $\cdot$ \bq@sub\ $\cdot$ \textbf{D\bq@diff}}\par\nopagebreak\relax}
  {\par{\small\color{teal}Answer: \textbf{\bq@ans}\ifx\bq@work\@empty\else\ (solved: \bq@work)\fi
  \ifx\bq@also\@empty\else\par\color{gray}Also in: \bq@also\fi}\par}
\define@key{bp}{id}{\def\bp@id{#1}}\define@key{bp}{source}{\def\bp@src{#1}}\define@key{bp}{questions}{\def\bp@qs{#1}}
\newenvironment{bankpassage}[1][]{\setkeys{bp}{#1}\par\needspace{4\baselineskip}\bigskip\noindent
  {\small\textbf{Passage \bp@id}\hfill\color{gray}[\bp@src; questions \bp@qs]}\par\nopagebreak\relax\itshape}{\par}
\makeatother
\newcommand{\bfig}[2]{\par\smallskip{\centering\includegraphics[width=#1pt,max width=0.62\linewidth,max height=0.32\textheight]{img/#2}\par}\smallskip}
\newcommand{\bfigi}[2]{\includegraphics[width=#1pt,max width=\linewidth,max height=0.32\textheight,valign=c]{img/#2}}
\newcommand{\bmissing}[1]{\textit{[figure missing in the source conversion: \detokenize{#1}]}}
\newenvironment{qoptions}{\begin{itemize}[leftmargin=1.8em,itemsep=1pt,topsep=2pt,label={}]}{\end{itemize}}
\newcommand{\opt}[2]{\item[(#1)] #2}
\providecommand{\fwstart}{}\providecommand{\fwend}{}\providecommand{\qfullsplit}{}
\newcolumntype{L}[1]{>{\raggedright\arraybackslash}p{#1}}
\newcolumntype{C}[1]{>{\centering\arraybackslash}p{#1}}
\newcolumntype{R}[1]{>{\raggedleft\arraybackslash}p{#1}}
\providecommand{\sectionhdr}[1]{}\providecommand{\subsectionhdr}[1]{}


\begin{document}
{\Large\bfseries Thermal Expansion \& Calorimetry --- Question Bank}\par\smallskip
{\small Physics chapter P13. Questions from all converted test papers, grouped by source series. Answers: printed key where it matches the paper, otherwise solved (answer\_src=solved, with working). Exact repeats across papers appear once, with also\_in listing the other copies.}\par\medskip\hrule

\section*{TOAS papers}
\begin{bankq}[id=TOAS11-P2-PHY-S1-Q06, source={TOAS - Test 11 (27 Jun 2022)}, paper=2, qno=6, type=multi, subtopic={Expansion}, answer={B,C}, answer_src=key, difficulty=4, diff_parts={K2 S2 I0 C0 F1}, diff_note={concepts: thermal expansion, elasticity of rods between walls | steps: free expansion of both rods; compression forces equal; strain and stress; force between rods}]
Two rods of different materials are placed between massive walls as shown in figure. The cross section of the rods is $A$, their moduli of elasticity are $E_1$ and $E_2$ respectively. If rods are heated by $t$ degrees, then (coefficients of linear expansion of material of rods are $\alpha_1$ and $\alpha_2$ respectively)
\bfig{138.5}{TOAS11_p2_phy_s1_q06_1.png}
\begin{qoptions}
\opt{A}{the strain produced in both rods is equal}
\opt{B}{the stress produced in both rods is equal}
\opt{C}{the force exerted by one rod on the other is $\dfrac{\alpha_1\ell_1+\alpha_2\ell_2}{\frac{\ell_1}{E_1}+\frac{\ell_2}{E_2}}At$}
\opt{D}{the force exerted by one rod on the other is $\dfrac{(E_1-E_2)^2A}{E_1+E_2}$}
\end{qoptions}
\end{bankq}
\begin{bankq}[id=TOAS11-P2-PHY-S3-Q05, source={TOAS - Test 11 (27 Jun 2022)}, paper=2, qno=5, type=numerical, subtopic={Calorimetry}, answer={12.00}, answer_src=key, difficulty=4, diff_parts={K1 S2 I0 C1 F1}, diff_note={concepts: calorimetry with variable specific heat | steps: heat lost by the liquid by integration; heat lost by the calorimeter; heat gained by the ice through melting and warming; solve m}]
A calorimeter of water equivalent 10 gm contains a liquid of mass 50 gm at $40^\circ$C. When `$m$' gram of ice at $-10^\circ$C is put into the calorimeter and the mixture is allowed to attain equilibrium, the final temperature was found to be 20 $^\circ$C. It is known that specific heat capacity of the liquid changes with temperature as $S=\left(1+\dfrac{\theta}{500}\right)$ Cal g$^{-1}$ $^\circ$C$^{-1}$ where $\theta$ is temperature in $^\circ$C. The specific heat capacity of ice, water and colorimeter remains constant and values are $S_{\text{ice}}=0.5$ Cal g$^{-1}$ $^\circ$C$^{-1}$, $S_{\text{water}}=1.0$ Cal g$^{-1}$ $^\circ$C$^{-1}$ and latent heat of fusion of ice is $L_f=80$ Cal g$^{-1}$. Assume no heat loss to the surrounding and calculate the value of `$m$' in gram.
\end{bankq}
\begin{bankq}[id=TOAS18-P2-PHY-S1-Q04, source={TOAS - Test 18 (30 Oct 2022)}, paper=2, qno=4, type=single, subtopic={Expansion}, answer={C}, answer_src=key, difficulty=3, diff_parts={K1 S2 I1 C0 F0}, diff_note={concepts: thermal expansion with a spacer | steps: capacitance dependence on area and gap; expansion of each part; condition}]
If surrounding temperature is slightly fluctuating as shown in figure. A parallel plate capacity with a spacer is design to operate at constant capacitance. If $\alpha_1$ \& $\alpha_2$ are coefficient of linear expansion of plate \& spacer respectively then. (Neglect area of spacer \& it is perfect electric insulator)
\bfig{82.0}{TOAS18_p2_phy_s1_q04_1.png}
\begin{qoptions}
\opt{A}{$\alpha_1=\alpha_2$}
\opt{B}{$\alpha_1=2\alpha_2$}
\opt{C}{$2\alpha_1=\alpha_2$}
\opt{D}{$2\alpha_1=3\alpha_2$}
\end{qoptions}
\end{bankq}
\begin{bankq}[id=TOAS19-P1-PHY-S1-Q07, source={TOAS - Test 19 (05 Feb 2023)}, paper=1, qno=7, type=multi, subtopic={Calorimetry}, answer={B,C}, answer_src=key, difficulty=3, diff_parts={K1 S2 I0 C0 F1}, diff_note={concepts: mixing of ice and water | steps: heat available from water; mass melted; remaining ice; effect of the calorimeter}]
20 gm ice is kept in a calorimeter at $0^\circ$C. Now 40 gm water at $30^\circ$C is poured gently into it. There is no leakage of heat. In equilibrium.
\begin{qoptions}
\opt{A}{Final temperature will be $0^\circ$C only if heat capacity of calorimeter is neglected.}
\opt{B}{Final temperature will be $0^\circ$C irrespective of value of heat capacity of calorimeter.}
\opt{C}{75\% of ice will melt.}
\opt{D}{Remaining ice will be more than 10 gm.}
\end{qoptions}
\end{bankq}

\section*{TOPS papers}
\begin{bankq}[id=TOPS04-P1-PHY-S1-Q06, source={TOPS - Test 04 (18 Dec 2022)}, paper=1, qno=6, type=multi, subtopic={Expansion}, answer={A,D}, answer_src=key, difficulty=6, diff_parts={K2 S2 I1 C1 F1}, diff_note={concepts: thermal expansion, angular momentum | steps: moment of inertia change; angular momentum conservation; new angular frequency}]
A uniform metallic circular disc of mass $M$ and radius $R$, mounted on frictionless bearings, is rotating an angular frequency $\omega$ about an axis passing through its centre and perpendicular to its plane. The temperature of the disc is then increased by $\Delta t$. If $\alpha$ is the coefficient of linear expansion of the metal,
\begin{qoptions}
\opt{A}{The moment of inertia increases by $MR^2\alpha\Delta t$.}
\opt{B}{The moment of inertia remains unchanged.}
\opt{C}{The angular frequency increases by $2\alpha\omega\Delta t$.}
\opt{D}{The angular frequency decreases by $2\alpha\omega\Delta t$.}
\end{qoptions}
\end{bankq}
\begin{bankq}[id=TOPS04-P2-PHY-S2-Q07, source={TOPS - Test 04 (18 Dec 2022)}, paper=2, qno=7, type=numerical, subtopic={Calorimetry}, answer={495.00}, answer_src=key, difficulty=7, diff_parts={K2 S2 I2 C1 F1}, diff_note={concepts: variable specific heat, calorimetry | steps: heat given by the container; heat taken by the ice; mass}]
An ice cube at $0^\circ$C of mass 0.1 kg is placed in an isolated container which is at $227^\circ$C. The specific heat capacity $c$ of the material of the container varies with temperature $T$ according to the empirical relation $C=A+BT$, where $A=100$ cal kg$^{-1}$ K$^{-1}$ and $B=2\times10^{-2}$ cal kg$^{-1}$ K$^{-2}$. If final temperature of the container is $27^\circ$C, find the mass of the container (in grams). The specific latent heat of fusion of ice is $8\times10^4$ cal kg$^{-1}$ and the specific heat capacity of water is $10^3$ cal kg$^{-1}$ ($^\circ$C)$^{-1}$.
\end{bankq}
\begin{bankq}[id=TOPS05-P1-PHY-S3-Q05, source={TOPS - Test 05 (14 Jan 2023)}, paper=1, qno=5, type=numerical, subtopic={Calorimetry}, answer={31.00}, answer_src=key, difficulty=5, diff_parts={K1 S2 I1 C1 F1}, diff_note={concepts: calorimetry | steps: heat given by the water; heat to warm the ice; heat of melting and warming; mass}]
How many grams of ice at $-14^\circ$C are needed to cool 200 g of water from $25^\circ$C to $10^\circ$C ? Take specific heat of ice $=0.5$ cal/g/$^\circ$C and latent heat of ice $=80$ cal/g.
\end{bankq}
\begin{bankq}[id=TOPS05-P1-PHY-S3-Q08, source={TOPS - Test 05 (14 Jan 2023)}, paper=1, qno=8, type=numerical, subtopic={Expansion}, answer={3.96}, answer_src=solved, working={Gap $=L\alpha\Delta T=12.0\times11\times10^{-6}\times30=3.96\times10^{-3}$ m $=3.96$ mm; the printed key shows 4.22, which would need $\alpha\approx11.7\times10^{-6}$}, difficulty=3, diff_parts={K1 S1 I0 C1 F1}, diff_note={concepts: thermal expansion | steps: expansion over the rise in temperature; gap}]
A railway track (made of iron) is laid in winter when the average temperature is 18$^\circ$C. The track consists of sections of 12.0 m placed one after the other. How much gap (in mm) should be left between two such sections so that there is no compression during summer when the maximum temperature goes to 48$^\circ$C ? Coefficient of linear expansion of iron $=11\times10^{-6}$ /$^\circ$C.
\end{bankq}
\begin{bankq}[id=TOPS05-P2-PHY-S3-Q03, source={TOPS - Test 05 (14 Jan 2023)}, paper=2, qno=3, type=numerical, subtopic={Specific and latent heat}, answer={2}, answer_src=key, difficulty=8, diff_parts={K2 S3 I2 C1 F1}, diff_note={concepts: variable specific heat, averages | steps: heat from 0 to t by integral; average heat capacity; heat capacity at t; difference; n}]
A 1kg substance of specific heat capacity is given as $C=(A+Bt^2)$, where $A$ and $B$ are constant and $t$ is temperature in $^\circ$C, is heated. Magnitude of difference in average heat capacity $\left(\dfrac{\Delta Q}{\Delta t}\right)$ from $0^\circ$C to $t^\circ$C and heat capacity at $t^\circ$C is given as $\dfrac{4BT^2}{3n}$, find $n$.
\end{bankq}
\begin{bankq}[id=TOPS05-P2-PHY-S3-Q08, source={TOPS - Test 05 (14 Jan 2023)}, paper=2, qno=8, type=numerical, subtopic={Expansion}, answer={6}, answer_src=key, difficulty=7, diff_parts={K2 S2 I2 C1 F1}, diff_note={concepts: volume expansion, liquid in a container | steps: expansion of the container volume; expansion of the liquid; condition for constant height; ratio}]
Figure shows a container of capacity 500 cc and cross-section area 20 cm$^2$, filled with 300 cc liquid at $20^\circ$C. Coefficient of volume expansion of liquid is $\gamma_L$ and coefficient of linear expansion of container is $\alpha_C$. Find $\dfrac{22\alpha_C}{\gamma_L}$ for which `$h$' remain constant at all temperature.
\bfig{58.2}{TOPS05_p2_phy_s3_q08_1.png}
\end{bankq}

\section*{TEAS papers}
\begin{bankq}[id=TEAS01-P1-PHY-S1-Q02, source={TEAS - Test 01 (16 Apr 2023)}, paper=1, qno=2, type=multi, subtopic={Calorimetry}, answer={A,B,D}, answer_src=key, difficulty=8, diff_parts={K2 S3 I2 C1 F1}, diff_note={concepts: calorimetry, graph reading | steps: heat balance for cooling water; the temperature when melting starts; ice mass from the graph; final state; check each}]
A calorimeter was filled with 10 g of water at a temperature of 70$^\circ$C. To cool it down, the researchers threw in calorimeter some crushed ice from a freezer at temperature $-10\,^\circ$C. Latent heat of ice $=330$ kJ/kg, specific heat capacity of water and ice are $c_w=4200$ J/(kg $^\circ$C) and $c_i=2100$ J/(kg $^\circ$C) respectively Graphs of dependences of the mass of each of substances from time are given on the image. Unfortunately, scientists have spilled coffee on the graph, and part of the experimental data was lost. Help researchers recover these data by choosing the correct answers. Neglect any heat loss.
\bfig{153.3}{TEAS01_p1_phy_s1_q02_1.png}
\begin{qoptions}
\opt{A}{The temperature of the water by the time the ice began to melt was 60$^\circ$C.}
\opt{B}{The final temperature in the calorimeter was 0$^\circ$C.}
\opt{C}{In steady state, 50\% was water and rest ice.}
\opt{D}{The mass of ice initially was 20 gm.}
\end{qoptions}
\end{bankq}
\end{document}
